\scanpage{47}
\mspage{23}
homomorphismes de groupes profinis associés, \deleted{regarder} \inserted{étudier} la situation analogue pour des ``modèles'' élémentaires \inserted{anabéliens}, voire des modèles $K(\pi,1)$ généraux. (On peut aussi regarder \unsure{de|des} tels modèles sur un corps \deleted{\unsure{de nombres}} \inserted{$K$ ext.} de type fini de $\mathbb Q$ -- mais passons pour le moment sur cette situation mixte, un peu bâtarde\ldots) Si $U,V$ sont des tels modèles, tout morphisme $\unsure{V|U}-U$ définit un morphisme de topos galoisiens sur $B_{\mathbb Q}$, $B_U\longrightarrow B_V$, \deleted{\unsure{de tels|dont le|dans le}} et si $\unsure{V|U}$ est \deleted{anabélien} élémentaire anabélien, cet \unsure{homph} est connu quand on connait seulement
\[
H_1(B_{\bar U},\mathbb Z_{\ell})-H_{\unsure{1|n}}(B_{\bar V},\mathbb Z_{\ell})
\]
-- ce qui est beaucoup moins que la classe d'iso. du hom. de $B_{\mathbb Q}$-topos. (En fait, \deleted{il \unsure{suffit}} \inserted{sans hyp} anabélienne sur $V$, dès que $V$ \unsure{se|est} plonge dans une VA, $f$ est connu quand on connait son action sur les types étales\ldots). Mais quels sont les hom $B_U-B_V$, ou $\mathcal E_U-\mathcal E_V$
\margin{46}
